% GENERATED FILE. Edit canonical lessons, then run npm run generate:books.
% canonical-source-hash: fnv1a64:787f872f204c31a7
% canonical-lessons: RU-C35-son, RU-C35-daughter, RU-C35-grandmother, RU-C35-husband, RU-C35-wife

\chapter{Son, Daughter, Grandmother, Husband, Wife}
\label{ch:ru-son-daughter-grandmother-husband-wife}
\hlchaptermodality{35}

\begin{chapteropening}
\textbf{By the end of this chapter:} \emph{I can name a son, a daughter, a grandmother, a husband and a wife.}

Complete the last lesson of chapter 35: i can name a son, a daughter, a grandmother, a husband and a wife.
\end{chapteropening}

\section[syn]{\ru{сын} (syn) — a son}

\label{lesson:RU-C35-son}



\begin{quote}
Before the new one: say the Russian for sadness, then the Russian for tiredness.
\end{quote}

\subsection*{You'll want to know: \ru{сын}}
\textbf{\ru{сын}} (\emph{syn}) — \textquotedblleft{}a son\textquotedblright{}.

\begin{grammarlens}[title={What you've built}]
One more word for son, daughter, grandmother, husband, wife.
\end{grammarlens}

\subsection*{Your turn}
\begin{itemize}
\raggedright
  \item \emph{Say it:} \emph{syn}
  \item \emph{Say it:} \emph{syn}, once more
  \item \emph{Say it:} say \emph{syn}
  \item \emph{Recall:} say the Russian for sadness, then the Russian for tiredness, then say \emph{syn} again
\end{itemize}

\begin{tcolorbox}[breakable,colback=teal!4,colframe=teal!35!black,title={Before you move on}]
What does \ru{сын} mean? (\textquotedblleft{}A son\textquotedblright{}.) Say it once more.
\end{tcolorbox}
\section[doch]{\ru{дочь} (doch) — a daughter}

\label{lesson:RU-C35-daughter}



\begin{quote}
Before the new one: say the Russian for tiredness, then the Russian for a son.
\end{quote}

\subsection*{You'll want to know: \ru{дочь}}
\textbf{\ru{дочь}} (\emph{doch}) — \textquotedblleft{}a daughter\textquotedblright{}.

\begin{grammarlens}[title={What you've built}]
One more word for son, daughter, grandmother, husband, wife.
\end{grammarlens}

\subsection*{Your turn}
\begin{itemize}
\raggedright
  \item \emph{Say it:} \emph{doch}
  \item \emph{Say it:} \emph{doch}, once more
  \item \emph{Say it:} say \emph{doch}
  \item \emph{Recall:} say the Russian for tiredness, then the Russian for a son, then say \emph{doch} again
\end{itemize}

\begin{tcolorbox}[breakable,colback=teal!4,colframe=teal!35!black,title={Before you move on}]
What does \ru{дочь} mean? (\textquotedblleft{}A daughter\textquotedblright{}.) Say it once more.
\end{tcolorbox}
\section[bábushka]{\ru{бабушка} (bábushka) — a grandmother}

\label{lesson:RU-C35-grandmother}



\begin{quote}
Before the new one: say the Russian for a son, then the Russian for a daughter.
\end{quote}

\subsection*{You'll want to know: \ru{бабушка}}
\textbf{\ru{бабушка}} (\emph{bábushka}) — \textquotedblleft{}a grandmother\textquotedblright{}.

\begin{grammarlens}[title={What you've built}]
One more word for son, daughter, grandmother, husband, wife.
\end{grammarlens}

\subsection*{Your turn}
\begin{itemize}
\raggedright
  \item \emph{Say it:} \emph{bábushka}
  \item \emph{Say it:} \emph{bábushka}, once more
  \item \emph{Say it:} say \emph{bábushka}
  \item \emph{Recall:} say the Russian for a son, then the Russian for a daughter, then say \emph{bábushka} again
\end{itemize}

\begin{tcolorbox}[breakable,colback=teal!4,colframe=teal!35!black,title={Before you move on}]
What does \ru{бабушка} mean? (\textquotedblleft{}A grandmother\textquotedblright{}.) Say it once more.
\end{tcolorbox}
\section[muzh]{\ru{муж} (muzh) — a husband}

\label{lesson:RU-C35-husband}



\begin{quote}
Before the new one: say the Russian for a daughter, then the Russian for a grandmother.
\end{quote}

\subsection*{You'll want to know: \ru{муж}}
\textbf{\ru{муж}} (\emph{muzh}) — \textquotedblleft{}a husband\textquotedblright{}.

\begin{grammarlens}[title={What you've built}]
One more word for son, daughter, grandmother, husband, wife.
\end{grammarlens}

\subsection*{Your turn}
\begin{itemize}
\raggedright
  \item \emph{Say it:} \emph{muzh}
  \item \emph{Say it:} \emph{muzh}, once more
  \item \emph{Say it:} say \emph{muzh}
  \item \emph{Recall:} say the Russian for a daughter, then the Russian for a grandmother, then say \emph{muzh} again
\end{itemize}

\begin{tcolorbox}[breakable,colback=teal!4,colframe=teal!35!black,title={Before you move on}]
What does \ru{муж} mean? (\textquotedblleft{}A husband\textquotedblright{}.) Say it once more.
\end{tcolorbox}
\section[zhená]{\ru{жена} (zhená) — a wife}

\label{lesson:RU-C35-wife}



\begin{quote}
Before the new one: say the Russian for a grandmother, then the Russian for a husband.
\end{quote}

\subsection*{You'll want to know: \ru{жена}}
\textbf{\ru{жена}} (\emph{zhená}) — \textquotedblleft{}a wife\textquotedblright{}.

\begin{grammarlens}[title={What you've built}]
One more word for son, daughter, grandmother, husband, wife.
\end{grammarlens}

\subsection*{Your turn}
\begin{itemize}
\raggedright
  \item \emph{Say it:} \emph{zhená}
  \item \emph{Say it:} \emph{zhená}, once more
  \item \emph{Say it:} say \emph{zhená}
  \item \emph{Recall:} say the Russian for a grandmother, then the Russian for a husband, then say \emph{zhená} again
\end{itemize}

\begin{tcolorbox}[breakable,colback=teal!4,colframe=teal!35!black,title={Before you move on}]
What does \ru{жена} mean? (\textquotedblleft{}A wife\textquotedblright{}.) Say it once more.
\end{tcolorbox}
